% !TeX root = ../../Book3.tex
% 纲要标题：### 19.3 完全交
% 纲要位置：计划纲要.md，第 2569 行。

\section{完全交}
\label{b3:sec:1903}

在一般环境环中，仅数方程和余维，不能保证每个方程在前面的商上仍为非零因子。在正则局部环中则有更强的判据：真理想的生成元数若等于其高度，这组生成元就是正则序列。完全交的内在定义要求，环的完备化能写成完备正则局部环的正则序列商。这样的方程组有 Koszul 消解，因而维数与深度可以同时计算。


% 纲要内容（仅作写作计划，尚非教材正文）：
%
% * 完全交局部环的完备化定义；用极大理想极小生成组的第一 Koszul 同调比较任意正则局部环境，证明完全交等价于当前核的最小生成元数等于高度，并区分极小与冗余方程组；实际正则环境中的局部化直接证明，一般完备化定义下的局部化注明所用平坦局部下降定理
% * Koszul 消解与 Tor 的变量交换；CM 模在正则局部环境中的投射维数等于支集余维；区分环境环有限消解与商环周期消解
% * 用两套 Koszul 消解计算同一 Tor，证明零维完全交的基座一维；平方零族给出 CM 而非完全交的反例
% * 正则蕴含完全交，完全交蕴含 CM；集中展示逆向反例及 CM 与所有 Serre 条件的等价关系
%

\subsection{正则环中正则序列定义的商}
\label{b3:subsec:1903:01}

设 \((S,\mathfrak n)\) 为正则局部环。给定 \(c\) 个方程后，若它们生成的真理想恰有高度 \(c\)，便能把这组方程补成 \(S\) 的完整参数系。\(S\) 是 CM 环，因而整组参数正则，原来的方程也就逐项正则。下面的\cref{b3:prop:given-complete-intersection}把这个判据与正则序列商的性质一起陈述。这里理想的高度仍取包含它的素理想高度的最小值；证明先计算商环的维数，以确定还需要补几个参数。

\begin{definition}\label{b3:def:local-complete-intersection}
设 \((R,\mathfrak m)\) 为交换 Noether 局部环，\(\widehat R\) 为其 \(\mathfrak m\) 进完备化。若存在完备正则局部环 \((S,\mathfrak n)\) 及 \(S\) 正则序列 \(f_1,\ldots,f_c\in\mathfrak n\)，使
\[
\widehat R\cong S/(f_1,\ldots,f_c),
\]
则称 \(R\) 为\term[wanquanjiaojubuhuan]{完全交局部环}[complete intersection local ring]。允许 \(c=0\)。对交换 Noether 环 \(A\)，若每个素理想 \(\mathfrak p\subseteq A\) 处的局部环 \(A_{\mathfrak p}\) 都是完全交局部环，则称 \(A\) 为\term[jubuwanquanjiaohuan]{局部完全交环}[locally complete intersection ring]。
\end{definition}

完全交性要求存在某个正则序列展示。在一个展示 \(R=S/I\) 中，\(I\) 的某组可能冗余的生成元不正则，不能仅据此否定完全交性。例如，重复写出同一个方程便会破坏正则性，却没有改变商环。不过，当 \(S\) 为正则局部环，且所取的是 \(I\) 的极小生成组时，情况不同：\cref{b3:thm:ci-presentation-independence}将证明，\(R\) 完全交当且仅当这组生成元正则；换用其他正则局部环境不能消除这一障碍。定义从 \(R\) 自身的完备化出发，因此不要求预先给出正则局部环到 \(R\) 的满射；完备局部环的环境构造与 Cohen 结构定理有关，见本卷附录。\cref{b3:prop:ci-localization}还将说明完全交局部环在其他素点处仍完全交，从而与一般环的局部完全交定义相容。

\ProseParagraphBreak
任意展示中的序列长度 \(c\) 也不能立刻看作 \(R\) 的内在数值。例如，对任意域 \(k\)，有
\[
k[\![x]\!]/(x^2)\cong k[\![x,T]\!]/(T,x^2).
\]
左边的正则序列长一，右边的正则序列长二，商环却相同；增加变量并增加消去该变量的关系，会同时增加环境维数与关系数。本章的局部完全交性是环自身在各素点的性质。代数几何中的局部完全交态射或正则嵌入则还涉及一个指定态射或嵌入，使用这些术语时须说明所指对象。

\begin{proposition}\label{b3:prop:given-complete-intersection}
设 \((S,\mathfrak n)\) 为交换 Noether 正则局部环，\(I=(f_1,\ldots,f_c)\subsetneq S\)，\(R=S/I\)。若 \(f_1,\ldots,f_c\) 为 \(S\) 正则序列，则 \(R\) 为完全交局部环，且 \(\dim R=\dim S-c\)。若 \(\operatorname{ht}I=c\)，则给定的 \(f_1,\ldots,f_c\) 必为 \(S\) 正则序列，因而 \(R\) 同样是维数为 \(\dim S-c\) 的完全交局部环。允许 \(c=0\)，此时 \(I=0\)。
\end{proposition}
\begin{proof}
先证明高度判据。记 \(d=\dim S\)。\cref{b3:thm:regular-local-characterizations}给出 \(\operatorname{depth}S=d\)，而\cref{b3:thm:depth-localization}与素链拼接给出，对每个 \(\mathfrak p\in\operatorname{Spec}S\)，
\[
d\le\operatorname{depth}S_{\mathfrak p}+\dim S/\mathfrak p
\le\operatorname{ht}\mathfrak p+\dim S/\mathfrak p\le d.
\]
中间一步使用 \(\operatorname{depth}S_{\mathfrak p}\le\dim S_{\mathfrak p}=\operatorname{ht}\mathfrak p\)。全部不等式于是成为等号，得到
\[
\operatorname{ht}\mathfrak p+\dim S/\mathfrak p=d.
\]
对包含 \(I\) 的素理想取商环维数的最大值，也就是对高度取最小值，便得
\(\dim S/I=d-\operatorname{ht}I=d-c\)。在 \(S/I\) 中选取 \(d-c\) 个参数 \(\bar g_1,\ldots,\bar g_{d-c}\)，并提升至 \(\mathfrak n\)。参数系的根条件给出
\[
\sqrt{(f_1,\ldots,f_c,g_1,\ldots,g_{d-c})}=\mathfrak n.
\]
因此这 \(d\) 个元素构成 \(S\) 的参数系。\(S\) 为 CM 环，\cref{b3:thm:cm-parameters}使整组参数正则，其前 \(c\) 项恰为给定方程。这也包括空序列的情形。

现在假定给定序列正则，不论它来自原假设还是刚证的高度判据。维数等式由\cref{b3:prop:regular-quotient-dimension}逐次应用得到。\cref{b3:thm:completion-exact}和\cref{b3:prop:completion-quotients}给出
\(\widehat R\cong\widehat S/(f_1,\ldots,f_c)\widehat S\)。\(S\to\widehat S\) 的平坦性保持各逐次商上的单射乘法；忠实平坦性则保证这些非零的商张量后仍非零。因此原序列在 \(\widehat S\) 上仍正则。再由\cref{b3:thm:regular-local-completion}，\(\widehat S\) 为完备正则局部环，所以上述展示符合完全交的定义。
\end{proof}

这个判据的关键是 CM 参数系定理。一般环境环缺少这项性质，不能只凭方程数与高度相等推出正则性。重复写一个方程也不会增加理想的高度，因而不能满足正则局部环中的这一判据。

\ProseParagraphBreak
设 \(k\) 为域。例如 \(k[\![x,y,z]\!]/(xy-z^2)\) 是二维完全交局部环：\(k[\![x,y,z]\!]\) 是三维正则局部整环，\(xy-z^2\) 是一个非零非单位。相比之下，\(k[\![x,y]\!]/(x^2,xy)\) 的两个给定方程不构成正则序列；在商去 \(x^2\) 后，非零的 \(x\) 被 \(xy\) 零化。

\subsection{完全交的 Koszul 消解}
\label{b3:subsec:1903:02}

\begin{proposition}\label{b3:prop:ci-koszul-resolution}\label{b3:prop:cm-pd-codimension}
设 \((S,\mathfrak n)\) 为交换 Noether 正则局部环，\(n=\dim S\)，\(k=S/\mathfrak n\) 为其剩余域。对任意非零有限生成的 \(S\) 模 \(M\)，令 \(d=\dim_SM\)，则
\[
M\;\text{为 CM 模}\quad\Longleftrightarrow\quad
\operatorname{pd}_SM=n-d.
\]
特别地，若 \(f_1,\ldots,f_c\in\mathfrak n\) 为 \(S\) 正则序列，\(R=S/(f_1,\ldots,f_c)\)，则 Koszul 复形 \(K(f_1,\ldots,f_c;S)\) 为 \(R\) 的极小自由 \(S\) 消解，其第 \(i\) 项的秩为 \(\binom ci\)。因而
\[
\dim_k\operatorname{Tor}_i^S(k,R)=\binom ci\quad(0\le i\le c),
\qquad\operatorname{pd}_S R=c.
\]
\end{proposition}
\begin{proof}
\cref{b3:thm:regular-local-homological-criterion}保证 \(M\) 的投射维数有限，\cref{b3:thm:regular-local-characterizations}给出 \(\operatorname{depth}S=n\)。因此可以应用\cref{b3:thm:auslander-buchsbaum}，得到
\[
\operatorname{pd}_SM=n-\operatorname{depth}_SM.
\]
右端等于 \(n-d\)，恰好意味着 \(\operatorname{depth}_SM=d\)，即 \(M\) 为 CM 模。

对所给的正则序列，\cref{b3:thm:koszul-regular-resolution}把其 Koszul 复形识别为 \(R\) 的自由消解。该复形第 \(i\) 项为 \(\bigwedge^iS^c\)。若 \(e_1,\ldots,e_c\) 为 \(S^c\) 的标准基，则这一项的基为
\[
e_{j_1}\wedge\cdots\wedge e_{j_i}
\qquad(1\le j_1<\cdots<j_i\le c),
\]
即从 \(c\) 个方程中选取 \(i\) 个所得到的 \(i\) 元子集，故其秩为 \(\binom ci\)。微分矩阵的各项是 \(0\) 或 \(\pm f_j\)，均属于 \(\mathfrak n\)，故此消解极小。与 \(k\) 张量后所有微分为零，首先得到
\[
\operatorname{Tor}_i^S(R,k)\cong k^{\binom ci}.
\]
再由\cref{b3:lem:tor-basic-properties}的自然同构
\(\operatorname{Tor}_i^S(R,k)\cong\operatorname{Tor}_i^S(k,R)\)，得到所列 Tor 维数。最高次 Tor 在 \(c\) 次非零，而自由消解长度为 \(c\)，故投射维数恰为 \(c\)。
\end{proof}

第一项等价把 CM 性翻译为环境环上的消解长度：若模的支集在正则局部环境中有余维 \(n-d\)，那么 CM 性恰好要求其投射维数达到这个值。正则序列商的 Koszul 消解实现了长度为 \(c\) 的这种情形。不过，CM 商环也可以有不是 Koszul 形状的极小消解，长度等于余维并不保证关系由正则序列生成。

\ProseParagraphBreak
以上的 \(\operatorname{pd}_SR=c\) 和 \(\operatorname{Tor}_i^S(k,R)\) 都以 \(S\) 为底环，所消解的是 \(S\) 模 \(R\)。改以 \(R\) 为底环时，\(R\) 本身自由，而剩余域 \(k\) 的投射维数可能无穷。例如，对域 \(k\) 和整数 \(n\ge2\)，后面的\cref{b3:prop:ci-cm-counterexample-families}给出 \(k[\![x]\!]/(x^n)\) 上的这种周期消解。完全交性提供正则环境环上的有限消解；商环剩余域消解的有限性则由\cref{b3:thm:regular-local-homological-criterion}刻画，对应于商环本身的正则性。

\subsection{完全交、Cohen–Macaulay 性与正则性的关系}
\label{b3:subsec:1903:03}

\begin{theorem}\label{b3:thm:complete-intersection-cm}
对于交换 Noether 局部环，有
\[
\text{正则}\ \overset{c=0}{\Longrightarrow}\ \text{完全交}
\ \overset{\text{正则序列商}}{\Longrightarrow}\ \text{Cohen--Macaulay}.
\]
\end{theorem}
\begin{proof}
正则环的完备化仍正则，取空序列即可得到第一项蕴含。若 \(R\) 为完全交局部环，则按定义 \(\widehat R=S/(f_1,\ldots,f_c)\)。\(S\) 为 CM 环，\cref{b3:thm:cm-regular-quotient}使该商为 CM 环，再由\cref{b3:thm:cm-completion}下降至 \(R\)。
\end{proof}

\(k[\![x]\!]/(x^2)\) 是完全交而非正则：其维数为零，极大理想却非零。局部环 \(B=k[\![x,y]\!]/(x^2,xy)\) 则不是 CM 环，因而不是完全交：\(B\) 的维数为一，而非零的 \(\bar x\) 被极大理想 \((\bar x,\bar y)\) 零化，所以基座非零，由\cref{b3:cor:depth-zero-socle}得到 \(\operatorname{depth}B=0\)。这里使用的是环自身的 CM 障碍，而非仅仅检查当前方程组。

\ProseParagraphBreak
要给出 CM 而非完全交的例子，可以在零维局部环上使用基座。设 \(k\) 为域，先看 \(C=k[\![x,y]\!]/(x^2,y^2)\)。它有基 \(1,x,y,xy\)，且对 \(a=a_0+a_1x+a_2y+a_3xy\)，条件 \(xa=ya=0\) 恰好要求 \(a_0=a_1=a_2=0\)。因此 \(\operatorname{Soc}(C)=k\cdot xy\) 为一维。相比之下，\(k[x,y]/(x,y)^2\) 的极大理想平方为零，其基座为二维的 \(k\cdot x\oplus k\cdot y\)。下面说明，前一种一维现象是所有零维完全交局部环都必须满足的条件。

\begin{lemma}\label{b3:lem:artinian-ci-socle}
设 \((R,\mathfrak m)\) 为零维交换 Noether 完全交局部环，\(k=R/\mathfrak m\) 为其剩余域。则其基座
\(\operatorname{Soc}(R)=(0:_R\mathfrak m)\) 满足 \(\dim_k\operatorname{Soc}(R)=1\)。
\end{lemma}
\begin{proof}
证明要对同一个 \(\operatorname{Tor}_c^S(k,R)\) 作两次计算：定义方程的 Koszul 消解把它算成 \(k\)，而环境环正则参数的 Koszul 消解把它识别为基座。

零维 Noether 局部环为 Artin 环，极大理想幂零，所以 \(R=\widehat R\)。写 \(R=S/(f_1,\ldots,f_c)\)，其中 \((S,\mathfrak n)\) 为完备正则局部环，\(f_1,\ldots,f_c\in\mathfrak n\) 为 \(S\) 正则序列。维数公式给出 \(c=\dim S\)。取 \(\mathfrak n\) 的一个极小生成组 \(x_1,\ldots,x_c\)，由\cref{b3:thm:regular-local-characterizations}，它是正则序列。它们在 \(R\) 中的像生成极大理想，即
\[
(x_1,\ldots,x_c)R=\mathfrak nR=\mathfrak m.
\]
且 \(S/\mathfrak n\cong R/\mathfrak m=k\)。

用 \(f\) 的 Koszul 消解计算，\cref{b3:prop:ci-koszul-resolution}给出 \(\operatorname{Tor}_c^S(k,R)\cong k\)。另一方面，用 \(x\) 的 Koszul 复形消解 \(k\)，再与 \(R\) 张量，其最高次同调为
\[
\begin{aligned}
H_c(K(x;S)\otimes_SR)
&=\{a\in R:x_1a=\cdots=x_ca=0\}\\
&=(0:_R\mathfrak m).
\end{aligned}
\]
这里最高项是 \(R\)，微分的各坐标为 \(\pm x_ja\)，且更高一项为零，故最高同调正是所写的共同零化子。两种计算得到同一个 Tor；其一致性由\cref{b3:lem:tor-basic-properties}的平衡性保证。因此基座的维数为一；\(c=0\) 时 \(R=k\)，结论同样成立。
\end{proof}

基座一维只是零维完全交局部环的必要条件，不能反过来作为完全交的充分判据；这个引理也不适用于前面一维的 \(k[\![x,y]\!]/(x^2,xy)\)，该环已由深度小于维数排除完全交性。

\begin{proposition}\label{b3:prop:ci-cm-counterexample-families}
设 \(k\) 为域。
\begin{enumerate}
\item 设 \(n\ge2\) 为整数，\(A=k[\![x]\!]/(x^n)\)。\(A\) 为零维完全交局部环，且交替使用乘 \(x\) 与乘 \(x^{n-1}\) 的复形
\[
\cdots\xrightarrow{x}A\xrightarrow{x^{n-1}}A
\xrightarrow{x}A\longrightarrow k\longrightarrow0
\]
为剩余域 \(k\) 的极小自由 \(A\) 消解。因此
\[
\operatorname{Tor}_i^A(k,k)\cong k\quad(i\ge0),
\qquad\operatorname{pd}_Ak=\infty.
\]
\item 设 \(e\ge2\) 为整数，
\[
B=k[u_1,\ldots,u_e]/(u_1,\ldots,u_e)^2,
\qquad\mathfrak m=(\bar u_1,\ldots,\bar u_e).
\]
则 \(B\) 为零维 CM 局部环，长度为 \(e+1\)，嵌入维数为 \(e\)，且
\[
\operatorname{Soc}(B)=\mathfrak m,
\qquad\dim_k\operatorname{Soc}(B)=e.
\]
因此 \(B\) 不是完全交局部环。
\end{enumerate}
\end{proposition}
\begin{proof}
对第一族，\(k[\![x]\!]\) 是一维正则局部整环，\(x^n\) 为非零非单位，故\cref{b3:prop:given-complete-intersection}说明 \(A\) 为零维完全交局部环。在 \(A\) 中，乘 \(x\) 的核为 \((x^{n-1})\)，乘 \(x^{n-1}\) 的核为 \((x)\)，而 \(A\to k\) 的核也为 \((x)\)，所以所给的列正合。两个乘法系数都在极大理想中，故消解极小；张量 \(k\) 后微分全部为零，得到每个次数的 Tor 都同构于 \(k\)。这些 Tor 在任意高次非零，因此 \(\operatorname{pd}_Ak=\infty\)。当 \(n=2\) 时，所有微分都是乘 \(x\)。

对第二族，\(B\) 的 \(k\) 向量空间基为 \(1,\bar u_1,\ldots,\bar u_e\)。每个元素唯一写成 \(a+v\)，其中 \(a\in k\)、\(v\in\mathfrak m\)；\(a\ne0\) 时由 \(v^2=0\) 得到逆元 \(a^{-1}-a^{-2}v\)，所以 \(\mathfrak m\) 是唯一极大理想。又 \(\mathfrak m^2=0\)，故 \(B\) 为零维局部环，自动为 CM 环。它的所有简单因子均为剩余域 \(k\)，因而长度等于其 \(k\) 维数，即 \(e+1\)；由 \(\mathfrak m/\mathfrak m^2=\mathfrak m\)，嵌入维数为 \(e\)。乘以 \(\bar u_i\) 时有 \(\bar u_i(a+v)=a\bar u_i\)，故同时被全部 \(\bar u_i\) 零化恰好意味着 \(a=0\)。因此基座正是 \(\mathfrak m\)，其维数为 \(e\ge2\)。\cref{b3:lem:artinian-ci-socle}遂排除了所有可能的完全交展示。
\end{proof}

\begin{center}
\begin{minipage}{\textwidth}
\makeatletter
\def\@captype{figure}
\makeatother
\centering
\[
\begin{tikzcd}[column sep=large]
\text{正则} \arrow[r, Rightarrow, "{c=0}"] &
\text{完全交} \arrow[r, Rightarrow, "{\text{正则序列商}}"] &
\text{Cohen--Macaulay}
\end{tikzcd}
\]
\[
\text{Cohen--Macaulay}
\quad\Longleftrightarrow\quad
\text{对每个整数}\;r\ge1\;\text{满足}\;(S_r).
\]
\begin{tabularx}{\textwidth}{@{}>{\centering\arraybackslash}X>{\centering\arraybackslash}X@{}}
完全交而非正则 & CM 而非完全交\\
\(k[\![x]\!]/(x^2)\) & \(k[x,y]/(x,y)^2\)
\end{tabularx}
\caption[正则、完全交与 Cohen--Macaulay 性]{正则、完全交与 Cohen--Macaulay 性。性质关系均限于交换 Noether 局部环；下方两个逆向反例中的 \(k\) 为任意域。}
\label{b3:fig:regular-ci-cm}
\end{minipage}
\end{center}

图中的第二行由\cref{b3:thm:cm-localization}与\cref{b3:def:serre-sk}得到。若 \(R\) 为 CM 局部环，则对每个素理想 \(\mathfrak p\)，局部化后有 \(\operatorname{depth}R_{\mathfrak p}=\dim R_{\mathfrak p}\)，自然满足任意 \((S_r)\)。反之，若 \(R\) 满足所有 \((S_r)\)，则在任一 \(\mathfrak p\) 处选取整数 \(r\ge\max\{1,\dim R_{\mathfrak p}\}\)，其条件便变成
\[
\operatorname{depth}R_{\mathfrak p}\ge\dim R_{\mathfrak p}.
\]
结合深度不超过维数，得到该点的 CM 等式；尤其在极大理想处得到 \(R\) 为 CM 环。这里的“所有 \(r\)”保证在每个点都能取到不小于局部维数的整数，因而恰好要求各局部环的深度达到其维数。

定义只要求存在一个合适的展示，而实际计算通常给出一个指定的正则局部环境。要使两者接上，需要比较不同环境中的关系数。极大理想的第一 Koszul 同调提供了这种比较：多添一个可以消去的环境参数，就同时多出一个零生成元的 Koszul 同调类；扣除这些类以后，剩下的数据只依赖商环。

\begin{theorem}\label{b3:thm:ci-presentation-independence}
设 \((S,\mathfrak n)\) 为交换 Noether 正则局部环，\(I\subsetneq S\) 为理想，\(R=S/I\)。记 \(\mu_S(I)=\dim_{S/\mathfrak n}I/\mathfrak nI\)，即 \(I\) 的最小生成元数。下列条件等价：
\begin{equivalences}
\item \(R\) 是完全交局部环；
\item \(I\) 可由一个 \(S\) 正则序列生成；
\item \(I\) 的任意极小生成组都是 \(S\) 正则序列；
\item \(\mu_S(I)=\operatorname{ht}_S I\)。
\end{equivalences}
\end{theorem}
\symidx{\(\mu_S(I)\)：局部环 \(S\) 中理想 \(I\) 的最小生成元数}
\begin{proof}
记 \(\mathfrak m=\mathfrak n/I\)、\(k=S/\mathfrak n=R/\mathfrak m\)、\(n=\dim S\)、\(e=\operatorname{edim}R\)、\(d=\dim R\)。取 \(\mathfrak m\) 的极小生成组 \(z=(z_1,\ldots,z_e)\)，并记
\[
h(R)=\dim_k H_1(z;R).
\]
\cref{b3:prop:koszul-signs-homotopy}说明 \(\mathfrak m\) 零化这个同调模；它又是有限生成模，因而上式是有限的 \(k\) 维数。这个数不依赖极小生成组：若 \(w\) 是另一组，则 \(w=Uz\) 对某个 \(e\times e\) 矩阵 \(U\) 成立。两组模 \(\mathfrak m^2\) 都是 \(\mathfrak m/\mathfrak m^2\) 的基，所以 \(U\) 模 \(\mathfrak m\) 可逆，\(\det U\) 为单位。\cref{b3:prop:koszul-generator-change}遂给出相应 Koszul 复形的同构。

取 \(\mathfrak n\) 的极小生成组 \(x_1,\ldots,x_n\)。由\cref{b3:thm:regular-local-characterizations}，它们是正则序列，故 \(K(x;S)\) 消解 \(k\)。它们在 \(R\) 中的像生成 \(\mathfrak m\)；从这些像中选出 \(e\) 个，使一次类组成 \(\mathfrak m/\mathfrak m^2\) 的基，重排以后记作 \(z_1,\ldots,z_e\)。Nakayama 引理保证它们生成整个 \(\mathfrak m\)。其余每个像都可写成这些 \(z_i\) 的 \(R\) 线性组合；从相应行减去这个组合，给出 \(R\) 上的可逆矩阵，把整组像变成
\[
(z_1,\ldots,z_e,\underbrace{0,\ldots,0}_{n-e\;\text{项}}).
\]
这里仅在张量后的复形 \(K(x;S)\otimes_SR\) 上换基，不需要把这个矩阵提升为 \(S\) 上的参数变换。由\cref{b3:prop:koszul-generator-change,b3:prop:koszul-tensor}，新增的零生成元贡献微分为零的外代数。总次数一分成 \(H_1(z;R)\otimes_R\bigwedge^0R^{n-e}\) 与 \(H_0(z;R)\otimes_R\bigwedge^1R^{n-e}\)；因 \(H_0(z;R)=R/\mathfrak m=k\)，得到
\[
\operatorname{Tor}_1^S(k,R)
\cong H_1(z;R)\oplus k^{n-e}.
\]
另一方面，对 \(0\to I\to S\to R\to0\) 使用\cref{b3:lem:tor-basic-properties}的 Tor 长正合列。\(S\) 自由，且 \(I\subseteq\mathfrak n\) 使 \(I/\mathfrak nI\to S/\mathfrak n\) 为零，故
\[
\operatorname{Tor}_1^S(k,R)\cong I/\mathfrak nI,
\qquad
\mu_S(I)=h(R)+n-e.
\]
前面高度判据证明中的素理想维数等式还给出
\(\operatorname{ht}_S I=n-d\)。因此，对任意这样的正则局部展示，都有
\[
\mu_S(I)=\operatorname{ht}_S I
\quad\Longleftrightarrow\quad h(R)=e-d.
\]
右边已经不含环境环 \(S\)。

这三个数在完备化下也不变。\cref{b3:thm:completion-dimension}给出 \(\dim\widehat R=d\)；\cref{b3:prop:completion-quotients}把 \(\mathfrak m/\mathfrak m^2\) 识别为 \(\mathfrak m\widehat R/(\mathfrak m\widehat R)^2\)，故嵌入维数仍为 \(e\)。有限自由 Koszul 复形经忠实平坦完备化后，其同调由\cref{b3:thm:completion-flat}变成 \(H_1(z;R)\otimes_R\widehat R\)。这个模被 \(\mathfrak m\) 零化，而完备化的剩余域仍为 \(k\)，所以它与原同调模同构。于是 \(h(\widehat R)=h(R)\)。

若第一项成立，按定义写 \(\widehat R=T/J\)，其中 \(T\) 为完备正则局部环，\(J\) 由长度 \(c\) 的正则序列生成。\cref{b3:prop:given-complete-intersection}给出 \(\operatorname{ht}_TJ=c\)；又有 \(\mu_T(J)\le c\)，而高度定理给出反向不等式，故 \(\mu_T(J)=\operatorname{ht}_TJ\)。把刚得到的关系数等式用于这个展示，得到
\[
h(\widehat R)=\operatorname{edim}\widehat R-\dim\widehat R.
\]
三个数的完备化不变性于是给 \(h(R)=e-d\)，再用于当前展示 \(R=S/I\)，得到第四项。

若第四项成立，任取 \(I\) 的极小生成组，其长度等于 \(\operatorname{ht}_SI\)，\cref{b3:prop:given-complete-intersection}保证这组元素正则，得到第三项。第三项蕴含第二项；第二项又由同一命题蕴含第一项。\(I=0\) 时它由空正则序列生成；\(e=0\) 时 \(z\) 为空，\(R=k\)、\(h(R)=0\)，相同的 Koszul 计算给出 \(\mu_S(I)=n\)，所以这些边界情形也包括在论证内。
\end{proof}

这个比较说明，判定强度取决于是否已经去掉冗余方程。若当前正则局部环境中一个极小生成组不正则，就可以否定环自身的完全交性；仅仅在某个非极小方程组中发现零因子，则还不能作这个结论。

完全交性还应当能从极大理想处传到其他素点。若环本身已经写成正则局部环的正则序列商，直接局部化即可；若只知道完备化有这样的展示，则还要从完备化上的素点下降回原环。这一步需要如下平坦局部下降定理：对交换 Noether 局部环之间的平坦局部同态 \(A\to B\)，若 \(B\) 是完全交局部环，则 \(A\) 也是完全交局部环。这项一般下降定理的证明不由平坦性检测正则序列直接得到；它通过正则局部覆盖及有限 Tor 维数比较两边的定义理想，使用分幂消解。其完整证明见 Stacks Project, Proposition 23.8.4 及 Lemma 23.7.6\footnote{分别见 \url{https://stacks.math.columbia.edu/tag/09Q2} 与 \url{https://stacks.math.columbia.edu/tag/09PX}。这里采用前一命题的下降方向；完整比较还涉及闭纤维的完全交性。本书不在此展开分幂消解与相应有限 Tor 维数判据的证明。}；下面使用这项输入证明局部化结论。

\begin{proposition}\label{b3:prop:ci-localization}
设 \((R,\mathfrak m)\) 为交换 Noether 完全交局部环。对每个 \(\mathfrak p\in\operatorname{Spec}R\)，\(R_{\mathfrak p}\) 都是完全交局部环。因此，完全交局部环也满足\cref{b3:def:local-complete-intersection}中局部完全交环的条件。
\end{proposition}
\begin{proof}
先处理具有实际正则局部展示的情形。若 \(R=S/(f_1,\ldots,f_c)\)，其中 \(S\) 为正则局部环、\(f\) 为 \(S\) 正则序列，取 \(\mathfrak p\) 在 \(S\) 中的原像 \(\mathfrak q\)。由\cref{b3:thm:regular-local-localization}，\(S_{\mathfrak q}\) 正则。局部化的正合性保持每个逐次商上的单射乘法，而全部 \(f_i\in\mathfrak q\) 保证末端商仍非零，因此这个序列在 \(S_{\mathfrak q}\) 上正则。由
\[
R_{\mathfrak p}\cong S_{\mathfrak q}/(f_1,\ldots,f_c)S_{\mathfrak q}
\]
及\cref{b3:prop:given-complete-intersection}，得到 \(R_{\mathfrak p}\) 完全交。这一论证也适用于由\cref{b3:thm:ci-presentation-independence}恢复正则序列核的指定正则局部展示。

一般情形中，定义给出 \(\widehat R=S/(f_1,\ldots,f_c)\)。\cref{b3:thm:completion-flat}保证 \(R\to\widehat R\) 忠实平坦，所以其素谱映射满射。取 \(\mathfrak q\in\operatorname{Spec}\widehat R\)，使 \(\mathfrak q\cap R=\mathfrak p\)。于是
\[
R_{\mathfrak p}\longrightarrow (\widehat R)_{\mathfrak q}
\]
是平坦局部同态：平坦性由局部化保持，目标极大理想的原像为 \(\mathfrak pR_{\mathfrak p}\)。闭纤维非零，\cref{b3:thm:faithfully-flat-criterion}还保证它忠实平坦。前一段已说明 \((\widehat R)_{\mathfrak q}\) 完全交；现在使用所述平坦局部下降定理，得到 \(R_{\mathfrak p}\) 完全交。
\end{proof}

这里没有把 \((R_{\mathfrak p})^{\wedge}\) 与 \((\widehat R)_{\mathfrak q}\) 视为同一个环；两者通过上面的平坦局部比较传递完全交性。一般地，也不能仅由 \(A\to B\) 平坦就把完全交性任意双向传递，向上保持还须对闭纤维加条件。
